\chapter{加法和减法}\label{ch:g3:addsub}

加法表示把数合在一起，减法表示去掉一部分，或
求出缺少的部分。本章训练两种竖式方法
——包括著名的进位和退位——以及
常常让竖式变得不必要的口算技巧。

\section{竖式加法}

\begin{method}[竖式加法]\label{met:g3:addsub:add}
\begin{enumerate}
  \item 把各数上下写齐，个位对个位，
        十位对十位（\emph{右对齐}）；
  \item 个位相加；若和超过 $9$，则写下和的个位数字，
        并将一个十\emph{进位}到下一列（写作小 $1$）；
  \item 继续对十位、百位、\dots 逐位相加，
        每次都要加上进位数。
\end{enumerate}
\end{method}

\begin{example}\label{ex:g3:addsub:add}
计算 $457 + 368$：
\[
\begin{array}{r}
{\scriptstyle 1}\ \ {\scriptstyle 1}\ \ \phantom{0} \\[-3pt]
4\,5\,7 \\
+\ 3\,6\,8 \\
\hline
8\,2\,5
\end{array}
\]
个位：$7 + 8 = 15$，写下 $5$，进 $1$。十位：
$5 + 6 + 1 = 12$，写下 $2$，进 $1$。百位：$4 + 3 + 1 = 8$。
结果为 $825$。
\end{example}

\begin{example}[口算加法]\label{ex:g3:addsub:mental}
在列竖式之前，先找一找好算的数：
\begin{itemize}
  \item \emph{凑十}：$47 + 8 = 47 + 3 + 5 = 50 + 5 = 55$；
  \item \emph{分步跳算}：$256 + 130 = 256 + 100 + 30 = 386$；
  \item \emph{调换顺序}：$17 + 59 + 3 = (17 + 3) + 59 = 79$。
\end{itemize}
\end{example}

\begin{omfigure}
\begin{tikzpicture}[scale=1.05]
  \draw[-Stealth] (-0.2,0) -- (10.4,0);
  \foreach \x/\l in {0.56/{256}, 5.56/{356}, 7.06/{386}}
    { \draw (\x,-0.08) -- (\x,0.08);
      \node[below, font=\small] at (\x,-0.1) {$\l$}; }
  \draw[-Stealth, omDef, thick] (0.56,0.25) .. controls (2.8,1.15) ..
    (5.56,0.25);
  \node[omDef, font=\small] at (3.05,1.05) {$+100$};
  \draw[-Stealth, omThm, thick] (5.56,0.25) .. controls (6.3,0.8) ..
    (7.06,0.25);
  \node[omThm, font=\small] at (6.3,0.85) {$+30$};
\end{tikzpicture}

{\small 用数轴上的跳跃表示口算加法：$256 + 130$ 分两次
跳完，先加百位，再加十位。}
\end{omfigure}

\section{竖式减法}

\begin{method}[竖式减法]\label{met:g3:addsub:sub}
\begin{enumerate}
  \item 把较大的数写在上面，右端对齐；
  \item 从个位开始，逐位相减；
  \item 若上方数字不够减，就\emph{退位}：从上一行的高一位借一个，
        它在当前列相当于十，
        计入当前列；
  \item \emph{验算}：差加上减数
        应等于上方的数。
\end{enumerate}
\end{method}

\begin{example}\label{ex:g3:addsub:sub}
计算 $603 - 147$：
\[
\begin{array}{r}
6\,0\,3 \\
-\ 1\,4\,7 \\
\hline
4\,5\,6
\end{array}
\]
个位：$3 - 7$ 不够减，需要连续退位：$603$ 变为
$5$ 个百、$9$ 个十、$13$ 个一。然后 $13 - 7 = 6$；
十位 $9 - 4 = 5$；百位 $5 - 1 = 4$。结果为 $456$。
验算：$456 + 147 = 603$。
\end{example}

\begin{omfigure}
\begin{tikzpicture}[scale=0.82]
  % before: 6 hundreds, 0 tens, 3 units
  \node[font=\small] at (2.5,2.25) {$603$ 退位前};
  \foreach \i in {0,...,5}
    \draw[omProp, fill=omProp!25] (0.62*\i,0.7) rectangle
      (0.62*\i+0.44,1.7);
  \node[below, font=\footnotesize] at (1.7,0.55) {$6$ 个百};
  \node[font=\footnotesize] at (4.35,1.2) {$0$ 个十};
  \foreach \i in {0,1,2} \fill[omDef] (5.4+0.4*\i,1.2) circle (0.13);
  \node[below, font=\footnotesize] at (5.8,0.55) {$3$ 个一};
  % after: 5 hundreds, 9 tens, 13 units
  \begin{scope}[yshift=-3.3cm]
  \node[font=\small] at (2.5,2.25) {$603$ 退位后};
  \foreach \i in {0,...,4}
    \draw[omProp, fill=omProp!25] (0.62*\i,0.7) rectangle
      (0.62*\i+0.44,1.7);
  \node[below, font=\footnotesize] at (1.4,0.55) {$5$ 个百};
  \foreach \i in {0,...,8}
    \draw[omThm, very thick] (3.5+0.22*\i,0.7) -- (3.5+0.22*\i,1.7);
  \node[below, font=\footnotesize] at (4.4,0.55) {$9$ 个十};
  \foreach \i in {0,...,6} \fill[omDef] (6.1+0.4*\i,1.45) circle (0.13);
  \foreach \i in {0,...,5} \fill[omDef] (6.1+0.4*\i,0.95) circle (0.13);
  \node[below, font=\footnotesize] at (7.3,0.55) {$13$ 个一};
  \end{scope}
\end{tikzpicture}

{\small 借位的图示 $603 - 147$：一百拆成十个十，
其中一个十再拆成十个一。
数量不变，只是换了新装——$6\,\vert\,0\,\vert\,3$ 变成
$5\,\vert\,9\,\vert\,13$，每一列现在都可以相减了。}
\end{omfigure}

\begin{example}[把减法看作求缺少的部分]\label{ex:g3:addsub:missing}
“莉拉有 $35$ 颗弹珠，汤姆有 $52$ 颗：汤姆比莉拉多多少颗？”
就是减法 $52 - 35$。心算时，从 $35$ \emph{往上}数：
$35 \to 40$ 是 $5$，再由 $40 \to 52$ 是 $12$，合计
$5 + 12 = 17$。往上数往往比借位减法更容易。
\end{example}

\begin{omfigure}
\begin{tikzpicture}[scale=1.05]
  \draw[-Stealth] (-0.2,0) -- (10.4,0);
  \foreach \x/\l in {1.0/{35}, 3.5/{40}, 9.5/{52}}
    { \draw (\x,-0.08) -- (\x,0.08);
      \node[below, font=\small] at (\x,-0.1) {$\l$}; }
  \draw[-Stealth, omDef, thick] (1.0,0.25) .. controls (2.2,0.9) ..
    (3.5,0.25);
  \node[omDef, font=\small] at (2.25,0.92) {$+5$};
  \draw[-Stealth, omThm, thick] (3.5,0.25) .. controls (6.5,1.2) ..
    (9.5,0.25);
  \node[omThm, font=\small] at (6.5,1.1) {$+12$};
\end{tikzpicture}

{\small 往上数计算 $52 - 35$：先数到好算的
整十数 $40$，再数到 $52$。答案是各次跳跃之和：
$5 + 12 = 17$。}
\end{omfigure}

\section{数量级}

\begin{method}[先估算再计算]\label{met:g3:addsub:estimate}
在做竖式计算之前，先把各数四舍五入成好算的数，
并算出一个\emph{估计值}。计算之后再比较：结果与估计值相差很大的，
说明可能算错了。
\end{method}

\begin{example}
对于 $487 + 312$：估计值为 $500 + 300 = 800$。竖式
结果 $799$ 很接近，是合理的。如果你算出了 $1\,099$
（数位没有对齐），估计值就会提醒你。
\end{example}

\section{练习}

\begin{exercise}[$\star$]\label{exo:g3:addsub:1}
用竖式计算：$346 + 253$；\quad $478 + 265$；\quad
$1\,536 + 2\,876$。
\end{exercise}

\begin{exercise}[$\star$]\label{exo:g3:addsub:2}
用竖式计算，再用加法验算：$587 - 234$；\quad
$805 - 348$；\quad $1\,000 - 463$。
\end{exercise}

\begin{exercise}[$\star$]\label{exo:g3:addsub:3}
用凑十法口算：$38 + 7$；\quad $56 + 9$；\quad
$45 + 28$（加 $30$，再减去 $2$）。
\end{exercise}

\begin{exercise}[$\star$]\label{exo:g3:addsub:4}
用跳跃法口算：$340 + 220$；\quad $675 + 210$；\quad
$512 - 300$；\quad $840 - 220$。
\end{exercise}

\begin{exercise}[$\star$]\label{exo:g3:addsub:5}
用往上数的方法计算 $63 - 47$（如
\cref{ex:g3:addsub:missing}），再用竖式验算。
\end{exercise}

\begin{exercise}[$\star$]\label{exo:g3:addsub:6}
先估算（四舍五入到百位），再精确计算：
$492 + 218$；\quad $913 - 388$。
\end{exercise}

\begin{exercise}[$\star$]\label{exo:g3:addsub:7}
一所学校有 $348$ 名女生，$296$ 名男生。一共有多少名学生？
列出算式，计算，并用一句话作答。
\end{exercise}

\begin{exercise}[$\star$]\label{exo:g3:addsub:8}
一本书有 $224$ 页，诺拉已经读了 $178$ 页，还剩
多少页没有读？
\end{exercise}

\begin{exercise}[$\star$]\label{exo:g3:addsub:9}
填出问号处的数字（逐列计算）：
\[
\begin{array}{r}
3\,?\,7 \\
+\ ?\,2\,? \\
\hline
8\,6\,2
\end{array}
\]
三个问号处的数字只有一种填法吗？请说明理由。
\end{exercise}

\begin{exercise}[$\star\star$]\label{exo:g3:addsub:10}
一家面包店上午卖出 $145$ 个羊角面包，
下午卖出 $89$ 个；打烊时还剩 $17$ 个没卖出。这一天
共烤了多少个羊角面包？（两步计算。）
\end{exercise}

\begin{exercise}[$\star\star$]\label{exo:g3:addsub:11}
西奥说：“算 $500 - 199$ 时，我先算 $500 - 200$，再
加 $1$。”他说得对吗？请两种方法都算一算，并说明他的窍门。
\end{exercise}
